\documentclass[12pt]{article}

\usepackage[spanish]{babel}
\usepackage[utf8]{inputenc}
\usepackage{longtable}
\usepackage{geometry}
\usepackage{tabularx}
\usepackage{array}
\geometry{margin=2.5cm}

\title{Backlog y Planificación}
\date{}

\begin{document}

\maketitle

\section{BACKLOG INICIAL}

\renewcommand{\arraystretch}{1.2}

\begin{longtable}{|p{2.0cm}|p{1.0cm}|p{5cm}|p{2cm}|p{4cm}|p{1cm}|}

\hline
\textbf{Épica} & \textbf{ID} & \textbf{Historia de Usuario} & \textbf{MoSCoW} & \textbf{Criterio de aceptación} & \textbf{Talla} \\
\hline

Gestión de eventos & HU-1 &
Como usuario, quiero ver un listado de eventos culturales y de ocio en Sevilla para descubrir actividades interesantes. &
Must &
Se muestra una lista con eventos, su imagen, nombre, categoría y fecha. &
M \\
\hline

Gestión de eventos & HU-2 &
Como usuario, quiero acceder a una ficha detallada de un evento para conocer su descripción, precio, ubicación y reseñas. &
Must &
Al pulsar un evento, se abre una vista con toda esta información. &
M \\
\hline

Gestión de eventos & HU-3 &
Como usuario registrado, quiero poder proponer eventos públicos y privados. &
Must &
El evento se envía a revisión de un moderador antes de publicarse. &
M \\
\hline

Moderación & HU-4 &
Como moderador, quiero aprobar o rechazar los eventos propuestos por los usuarios. &
Should &
El moderador puede ver la lista de eventos pendientes de aprobación y editarlos. &
M \\
\hline

Moderación & HU-5 &
Como usuario, quiero que se me notifique cuando se aprueba o se rechaza un evento propuesto por mí. &
Should &
Las notificaciones llegan a una pestaña de notificaciones. &
S \\
\hline

Mapas y geolocalización & HU-6 &
Como usuario, quiero ver los eventos en un mapa para ubicar los más cercanos y abrir la localización en la app de mapas. &
Must &
El mapa muestra los eventos y permite abrir los detalles del evento donde habrá un mapa que se puede abrir en el buscador. &
L \\
\hline

Mapas y geolocalización & HU-7 &
Como usuario, quiero recibir notificaciones sobre eventos cercanos a mi ubicación. &
Should &
El sistema envía una notificación en función de la localización. &
M \\
\hline

Interacción social & HU-8 &
Como usuario, quiero subir fotos y escribir mensajes en el chat de eventos para compartir mi experiencia y poder hablar de forma privada con personas específicas. &
Could &
La app permite escribir mensajes y tomar y subir fotos que quedarán guardados en el chat. &
L \\
\hline

Interacción social & HU-9 &
Como usuario, quiero escribir reseñas y valoraciones sobre los eventos. &
Should &
Se puede dejar puntuación (1-5 estrellas) y comentario. &
M \\
\hline

Autentica
-ción y perfil & HU-10 &
Como usuario, quiero registrarme, iniciar sesión y editar perfil. &
Must &
Se implementa Firebase Auth. &
M \\
\hline

Recomenda
-ciones & HU-11 &
Como usuario, quiero recibir recomendaciones de eventos y rutas de eventos según mis intereses. &
Should &
El sistema sugiere eventos y rutas según la interacción del usuario en la app (eventos vistos, eventos apuntados...). &
L \\
\hline

Administra
-ción del sistema & HU-12 &
Como administrador, quiero gestionar usuarios y categorías. &
Must &
Panel de administrador. &
S \\
\hline

Scraping & HU-13 &
Como usuario, quiero que el listado de eventos se actualice periódicamente. &
Must &
Scrapeo de diferentes páginas de eventos de sevilla. &
L \\
\hline

Búsqueda de eventos & HU-14 &
Como usuario, quiero buscar eventos en base a un radio de distancia, por categorías, por título/descripción y por precio. &
Must &
En el listado de eventos habrá filtros de búsqueda para todos los parámetros. &
S \\
\hline

Unión a eventos & HU-15 &
Como usuario, quiero poder apuntarme a eventos. &
Should &
En los eventos que son públicos habrá un botón en los detalles del evento para alistarse. Si el evento es privado habrá que unirse mediante un enlace que se muestra al crear el evento. &
M \\
\hline

Moderación de imágenes & HU-16 &
Como moderador, quiero comprobar que el contenido que se sube no es inadecuado &
Could &
Al subirse una imagen habrá un modelo que compruebe que el contenido es adecuado y si no lo es la elimina. &
L \\
\hline

Generación de carteles con IA & HU-17 &
Como administrador, quiero mostrar un cartel atractivo de portada del evento. &
Could &
Una IA generativa creará una imagen en base a la descripción del evento y se mostrará en la portada. &
M \\
\hline

Gestión de calendarios & HU-18 &
Como usuario, quiero tener un calendario con los eventos a los que me aliste. &
Could &
Se mostrará un calendario con todos los eventos alistados. &
S \\
\hline

\end{longtable}

\section{PLANIFICACIÓN (EDT) INICIAL}

\subsection{Nivel 1: Proyecto ``Sevillaneando''}

\begin{enumerate}

\item \textbf{Fase 1 – Análisis y planificación}

\begin{enumerate}

\item Revisión de requisitos y alcance

\item Elaboración del backlog inicial

\item Planificación temporal y de recursos (cronograma inicial)

\item Creación del modelo conceptual

\end{enumerate}

\item \textbf{Fase 2 – Diseño}

\begin{enumerate}

\item Diseño de arquitectura técnica

\item Diseño de la interfaz (mockups, navegación, usabilidad)

\item Diseño de la base de datos (PostgreSQL + PostGIS)

\item Diseño de la API REST (endpoints principales)

\item Planificación de la gestión de usuarios y permisos

\end{enumerate}

\item \textbf{Fase 3 – Desarrollo paralelo (Frontend y Backend)}

\textbf{3A: Desarrollo del frontend (React Native)}

\begin{enumerate}

\item Configuración del entorno móvil (dependencias, entornos, etc)

\item Implementación del listado y detalle de eventos

\item Creación y moderación de eventos, categorías y usuarios

\item Integración del mapa (OpenStreetMaps)

\item Integración de búsqueda

\item Sistema de notificaciones

\item Integración de chat

\item Módulo de subida de fotos

\item Implementación de reseñas y valoraciones

\item Personalización de recomendaciones

\end{enumerate}

\textbf{3B – Desarrollo del backend (Node.js + NestJS)}

\begin{enumerate}

\item Configuración del servidor, DB e integración de autenticación.

\item Implementación de endpoints de eventos, usuarios y categorías. (CRUD, lógica de eventos privados y sistema de moderación)

\item Sistema de Notificaciones

\item Infraestructura de Chat: WebSockets para chat.

\item Integración de imágenes.

\item Motor de Scraping: Implementación de scrapers para captura periódica de eventos externos

\item Algoritmo de Recomendaciones: Lógica de filtrado colaborativo o basado en etiquetas para rutas e intereses.

\end{enumerate}

\item \textbf{Fase 4 – Pruebas y validación}

\begin{enumerate}

\item Pruebas unitarias (frontend y backend)

\item Pruebas de integración (API + app)

\item Pruebas de usabilidad con usuarios reales

\item Corrección de errores y mejoras

\end{enumerate}

\item \textbf{Fase 5 – Ampliaciones y Extras: Esta fase se ejecutará únicamente si el cronograma de las fases anteriores lo permite}

\begin{enumerate}

\item Moderación Inteligente

\item Generación de Carteles con IA

\item Gestión de Calendarios

\end{enumerate}

\item \textbf{Fase 6 – Despliegue y documentación}

\begin{enumerate}

\item Despliegue del backend (servidor o nube)

\item Publicación de la app (Google Play / TestFlight)

\item Elaboración de la memoria técnica y manual de usuario

\end{enumerate}

\end{enumerate}

\end{document}